\chapter{Происхождение цветосохраняющего квадратичного селектора}
\label{ch:v7-color-preserving-quadratic-selector-origin-gate}

\section{Неожиданное совпадение двух проекторов}

Виртуальный цветной мост не запускает лёгкий сектор из нуля. Требуется
квадратичный проектор, который делает неустойчивыми только бесцветные
целевые рёбра. Полный gauge-подъём уже классифицировал одиннадцать блоков.
Четыре из них являются точными gauge-синглетами:
\begin{equation}
 E_0=\{L_LY_R,\ X_LX_R,\ X_Le_R,\ Y_LY_R\}.
 \label{eq:v7-color-selector-singlet-support}
\end{equation}
Но это в точности опора ранее выведенного изотипического проектора:
\begin{equation}
 \boxed{P_G=P_I,\qquad \rank P_G=4.}
 \label{eq:v7-color-selector-projector-coincidence}
\end{equation}
Здесь $P_I$ был получен условием эквивалентности калибровочных
представлений концов ребра, тогда как $P_G$ определяется независимо как
проектор на инвариантные векторы полного gauge-действия. Совпадение не
использует метки ``старое'' и ``новое''.

\section{Проектор как ядро Casimir}

Пусть на полном рёберном носителе
\begin{equation}
 \mathcal C_G=\alpha_1Y^2+\alpha_2C_2(SU(2))
 +\alpha_3C_2(SU(3)),\qquad \alpha_i>0.
 \label{eq:v7-color-selector-positive-casimir}
\end{equation}
Каждое слагаемое неотрицательно, поэтому ядро суммы не зависит от
положительных относительных нормировок:
\begin{equation}
 \ker\mathcal C_G=(\mathcal E_{\rm new})^{G_{SM}},
 \qquad
 P_G=\mathbf1_{\{0\}}(\mathcal C_G)
 =\lim_{t\to\infty}e^{-t\mathcal C_G}.
 \label{eq:v7-color-selector-kernel-projector}
\end{equation}
Следовательно, неизвестные отношения gauge-couplings не создают порога.
В обычной GUT-нормировке
\begin{equation}
 \mathcal C_G=\frac35Y^2+C_2(SU(2))+C_2(SU(3))
 \label{eq:v7-color-selector-gut-casimir}
\end{equation}
спектр на проверяемых типах равен
\begin{equation}
 \Spec\mathcal C_G=\left\{0,\frac9{10},\frac85,3\right\}.
 \label{eq:v7-color-selector-casimir-spectrum}
\end{equation}
Поэтому тот же проектор имеет конечную полиномиальную запись
\begin{equation}
 P_G=\left(I-\frac{\mathcal C_G}{9/10}\right)
 \left(I-\frac{\mathcal C_G}{8/5}\right)
 \left(I-\frac{\mathcal C_G}{3}\right).
 \label{eq:v7-color-selector-polynomial-projector}
\end{equation}
Фундаментальной является спектральная проекция на нуль; полином служит
точным сертификатом на данном конечном спектре.

\section{Новая физическая декомпозиция одиннадцати рёбер}

Циклический проектор $P_C$ и gauge-синглетный проектор $P_G=P_I$
коммутируют. Они дают каноническое ортогональное разбиение
\begin{equation}
 11=2_{\rm conn}+2_{\rm virtual}+2_{\rm mass}+5_{\rm forbidden}.
 \label{eq:v7-color-selector-edge-partition}
\end{equation}
Его части равны
\begin{align}
 E_{\rm conn}&=E_C\cap E_0=\{L_LY_R,X_Le_R\},\nonumber\\
 E_{\rm virtual}&=E_C\setminus E_0=\{Q_LY_R,X_Lu_R\},\nonumber\\
 E_{\rm mass}&=E_0\setminus E_C=\{X_LX_R,Y_LY_R\}.
 \label{eq:v7-color-selector-four-sector-support}
\end{align}
Таким образом, два прежде опасных цветных ребра получают точное физическое
чтение виртуальных мостов, а четыре бесцветных рёбра состоят из двух
смешанных коннекторов и двух вектороподобных масс.

\section{Hodge-запуск без нарушения gauge-группы}

Повторим уже выведенную двухцветную Hodge-конструкцию, заменяя прежний
проектор объединения на $P_G$. Производная градуировка равна
\begin{equation}
 \Gamma_G=I-2P_G.
 \label{eq:v7-color-selector-grading}
\end{equation}
Одна норма отображения момента даёт
\begin{equation}
 \mathcal S_G(z)=
 \sum_{e\in E_0}(|z_e|^2-\mu^2)^2
 +\sum_{e\notin E_0}(|z_e|^4+2\mu^2|z_e|^2).
 \label{eq:v7-color-selector-hodge-potential}
\end{equation}
Поэтому вакуум имеет опору
\begin{equation}
 |z_e|=\mu\quad(e\in E_0),\qquad
 z_e=0\quad(e\notin E_0).
 \label{eq:v7-color-selector-vacuum}
\end{equation}
Все ненулевые поля являются полными gauge-синглетами. В частности,
$SU(3)_c\times SU(2)_L\times U(1)_Y$ не нарушается этим внутренним
вектороподобным вакуумом.

На редуцированном пространстве меток сигнатуры равны
\begin{equation}
 (n_-,n_0,n_+)_{0}^{\rm red}=(8,0,14),\qquad
 (n_-,n_0,n_+)_{*}^{\rm red}=(0,4,18).
 \label{eq:v7-color-selector-reduced-signatures}
\end{equation}
После раскрытия минимальных полных multiplet-размерностей получаем
\begin{equation}
 (n_-,n_0,n_+)_{0}^{\rm full}=(8,0,34),\qquad
 (n_-,n_0,n_+)_{*}^{\rm full}=(0,4,38).
 \label{eq:v7-color-selector-full-signatures}
\end{equation}

\section{Сохранение лёгкого хирального лептонного сектора}

На одно поколение два синглетных массовых блока образуют отображение
$\mathbb C_L\to\mathbb C_R^2$, а два дублетных --- отображение
$\mathbb C_L^2\to\mathbb C_R$. При ненулевых амплитудах оба имеют ранг
один, поэтому
\begin{equation}
 \dim\ker M_{\rm doublet}=1,
 \qquad
 \dim\ker M_{\rm singlet}^{*}=1.
 \label{eq:v7-color-selector-light-kernels}
\end{equation}
При трёх семействах остаются три лёгких левых дублета и три лёгких правых
заряженных синглета. Тем самым сохраняется хиральный лептонный индекс
$H_{15}$, тогда как вектороподобное дополнение становится тяжёлым.

\begin{theorem}[Двойное происхождение физического вакуумного проектора]
На полном одиннадцатирёберном носителе проектор на ядро положительного
gauge-Casimir в точности совпадает с ранее независимо выведенным
изотипическим проектором $P_I$. Он выбирает четыре и только четыре
gauge-синглетных ребра. Производная Hodge-градуировка запускает их,
оставляет два цветных циклических моста и пять нежелательных блоков
массивными, сохраняет полную gauge-группу и один лёгкий хиральный
лептонный пакет на поколение.
\end{theorem}

\begin{nogocriterion}[Граница положительного прохода]
Совпадение $P_G=P_I$ закрывает квадратичный селектор, но ещё не фиксирует
коэффициент виртуального determinant-взаимодействия, общий масштаб $\mu$,
семейные ориентации и полное Real-бимодульное вложение. Нельзя возвращать
цветные рёбра в ненулевой фундаментальный вакуум или выбирать порог
Casimir вручную: используется только точное ядро.
\end{nogocriterion}

\begin{protocol}
\begin{enumerate}
 \item gauge-синглетных новых рёбер: \textbf{$4$};
 \item ранг $P_G$: \textbf{$4$};
 \item совпадение $P_G=P_I$: \textbf{точное};
 \item зависимость ядра от положительных gauge-весов: \textbf{нет};
 \item ручной Casimir-порог: \textbf{нет};
 \item бесцветных циклических коннекторов: \textbf{$2$};
 \item цветных виртуальных мостов: \textbf{$2$};
 \item вектороподобных массовых рёбер: \textbf{$2$};
 \item запрещённых рёбер: \textbf{$5$};
 \item полная сигнатура нуля: \textbf{$(8,0,34)$};
 \item полная сигнатура вакуума: \textbf{$(0,4,38)$};
 \item нарушение $SU(3)_c$: \textbf{нет};
 \item нарушение электрослабой gauge-группы: \textbf{нет};
 \item лёгкий левый дублет на поколение: \textbf{$1$};
 \item лёгкий правый синглет на поколение: \textbf{$1$};
 \item цветосохраняющий квадратичный селектор: \textbf{получен};
 \item следующий гейт: \textbf{совместный гессиан singlet-вакуума и
 виртуального цикла}.
\end{enumerate}
\end{protocol}

Машинный сертификат сохранён в
\path{s2t_v7_color_preserving_quadratic_selector_origin_gate_results.json}.

\begin{thebibliography}{9}
\bibitem{v7-color-selector-krajewski}
T.~Krajewski,
\emph{Classification of Finite Spectral Triples},
Journal of Geometry and Physics 28 (1998), 1--30;
arXiv:hep-th/9701081.
\bibitem{v7-color-selector-jureit}
J.-H.~Jureit, T.~Schuecker, C.~A.~Stephan,
\emph{On a Classification of Irreducible Almost-Commutative Geometries V},
Journal of Mathematical Physics 50 (2009), 072303; arXiv:0901.3214.
\bibitem{v7-color-selector-spectral-action}
A.~H.~Chamseddine, A.~Connes,
\emph{The Spectral Action Principle},
Communications in Mathematical Physics 186 (1997), 731--750;
arXiv:hep-th/9606001.
\end{thebibliography}